\chapter{Functional Programming in Lattice}\label{ch:functional-programming}

\index{functional programming}\index{pipe()}\index{compose()}\index{identity()}

Lattice is not a purely functional language---it has mutable variables, side effects, and imperative loops.
But it provides a rich set of functional tools that, when used well, can make your code more concise, more composable, and easier to reason about.

In \cref{ch:lazy-iterators}, we built pipelines of iterator transformations.
This chapter takes those ideas further.
We will explore Lattice's built-in functional primitives---\lstinline|pipe()|, \lstinline|compose()|, and \lstinline|identity()|---then dive into the \lstinline|fn| standard library's currying, partial application, and collection utilities.
By the end, you will know when functional style shines in Lattice and when to reach for something else.

%% ───────────────────────────────────────────────────────────
\section{pipe(), compose(), and identity()}\label{sec:pipe-compose-identity}
\index{pipe()}\index{compose()}\index{identity()}

These three built-in functions form the foundation of functional composition in Lattice.
They are implemented directly in the runtime (\filename{src/runtime.c} and \filename{src/eval.c}), so they are always available without any imports.

\subsection{pipe() --- Threading Values Through Functions}\label{subsec:pipe}
\index{pipe()!usage}

The \lstinline|pipe()| function takes a value and one or more functions, then threads the value through each function left to right:

\begin{lstlisting}[language=Lattice, caption={Basic pipe usage}]
let result = pipe(5,
    |x| { x * 2 },
    |x| { x + 1 },
    |x| { x * x }
)
print(result)  // 121  (5 -> 10 -> 11 -> 121)
\end{lstlisting}

Each function receives the output of the previous one.
The first argument is the initial value; every subsequent argument must be a closure.

\begin{defnbox}{pipe()}
\lstinline|pipe(val, f1, f2, ..., fn)| computes \lstinline|fn(...f2(f1(val)))|.
It evaluates left to right: the value flows \emph{forward} through the chain, like water through a series of pipes.
\end{defnbox}

Without \lstinline|pipe()|, you would write deeply nested function calls:

\begin{lstlisting}[language=Lattice, caption={Nested calls vs.\ pipe}]
// Without pipe: read inside-out
let result = square(increment(double(5)))

// With pipe: read top-to-bottom
let result = pipe(5,
    |x| { x * 2 },
    |x| { x + 1 },
    |x| { x * x }
)
\end{lstlisting}

The pipe version reads in the order that operations happen---no mental stack unwinding required.

Here is a more realistic example, processing a string:

\begin{lstlisting}[language=Lattice, caption={String processing with pipe}]
let raw_input = "  Hello, Lattice World!  "

let processed = pipe(raw_input,
    |s| { s.trim() },
    |s| { s.lower() },
    |s| { s.replace(",", "") },
    |s| { s.split(" ") }
)

print(processed)  // [hello, lattice, world!]
\end{lstlisting}

\begin{tipbox}{pipe() with Named Functions}
You are not limited to anonymous closures.
Named functions work too:
\begin{lstlisting}[language=Lattice]
fn double(x: Int) -> Int { return x * 2 }
fn increment(x: Int) -> Int { return x + 1 }

let result = pipe(5, double, increment)
print(result)  // 11
\end{lstlisting}
\end{tipbox}

\subsection{compose() --- Building New Functions from Old}\label{subsec:compose}
\index{compose()!usage}

While \lstinline|pipe()| immediately applies a chain of functions to a value, \lstinline|compose()| builds a \emph{new function} from two existing ones---without calling either:

\begin{lstlisting}[language=Lattice, caption={Composing functions}]
let double = |x| { x * 2 }
let increment = |x| { x + 1 }

// compose(f, g) returns a new function: |x| { f(g(x)) }
let double_then_increment = compose(increment, double)

print(double_then_increment(5))   // 11  (double(5)=10, increment(10)=11)
print(double_then_increment(10))  // 21  (double(10)=20, increment(20)=21)
\end{lstlisting}

\begin{warnbox}{Argument Order}
\lstinline|compose(f, g)| applies \lstinline|g| first, then \lstinline|f|.
In mathematical notation: \lstinline|compose(f, g)(x) = f(g(x))|.
This is right-to-left, like function composition in mathematics.
If you find this confusing, use \lstinline|pipe()| instead---it flows left-to-right.
\end{warnbox}

The value of \lstinline|compose()| is that it creates a reusable function you can pass around, store in variables, or use as a callback:

\begin{lstlisting}[language=Lattice, caption={Composed functions as callbacks}]
let to_celsius = |f| { (f - 32) * 5 / 9 }
let round_temp = |c| { to_int(c) }
let format_temp = |c| { to_string(c) + " C" }

// Build a single conversion function
let convert = compose(format_temp, compose(round_temp, to_celsius))

let temps_f = [98.6, 212.0, 32.0, 72.5]
let temps_c = iter(temps_f).map(convert).collect()
print(temps_c)  // [37 C, 100 C, 0 C, 22 C]
\end{lstlisting}

Under the hood, \lstinline|compose()| creates a new closure that captures both \lstinline|f| and \lstinline|g| as upvalues.
In the tree-walk evaluator (\filename{src/eval.c}), it builds an AST node for \lstinline|f(g(x))| with a fresh environment.
In the stack VM (\filename{src/runtime.c}), it generates bytecode: load \lstinline|g|, call with \lstinline|x|, load \lstinline|f|, call with the result.
Either way, the composed function behaves like any other closure.

\subsection{identity() --- The Do-Nothing Function}\label{subsec:identity}
\index{identity()!usage}

The \lstinline|identity()| function returns its argument unchanged:

\begin{lstlisting}[language=Lattice, caption={The identity function}]
print(identity(42))        // 42
print(identity("hello"))   // hello
print(identity([1, 2, 3])) // [1, 2, 3]
\end{lstlisting}

At first glance, this seems useless.
Why would you want a function that does nothing?
The answer is that \lstinline|identity()| serves as a \emph{default} or \emph{neutral element} in functional pipelines:

\begin{lstlisting}[language=Lattice, caption={identity() as a default transform}]
fn process_data(values: Array, transform: Fn) -> Array {
    return iter(values).map(transform).collect()
}

// Sometimes we want a transform, sometimes we don't
let raw = [1, 2, 3, 4, 5]

let doubled = process_data(raw, |x| { x * 2 })
print(doubled)  // [2, 4, 6, 8, 10]

let unchanged = process_data(raw, identity)
print(unchanged)  // [1, 2, 3, 4, 5]
\end{lstlisting}

Rather than adding special-case logic for ``no transformation,'' you pass \lstinline|identity| as the transform.
The code stays uniform.

Another use: \lstinline|identity| is the neutral element for \lstinline|compose()|.
Composing any function with \lstinline|identity| gives back the original function:

\begin{lstlisting}[language=Lattice, caption={identity() as composition neutral}]
let double = |x| { x * 2 }

let same = compose(double, identity)
print(same(5))  // 10
\end{lstlisting}

%% ───────────────────────────────────────────────────────────
\section{The fn Standard Library Module}\label{sec:fn-stdlib}
\index{fn standard library}\index{standard library!fn}

The \lstinline|fn| module (\filename{lib/fn.lat}) bundles a collection of functional programming utilities into a single import.
We encountered its lazy sequence functions in \cref{ch:lazy-iterators}; now let's explore the rest.

\begin{lstlisting}[language=Lattice, caption={Importing the fn module}]
import "lib/fn" as F
\end{lstlisting}

\subsection{Function Composition Utilities}\label{subsec:fn-composition}
\index{fn standard library!comp}\index{fn standard library!flip}\index{fn standard library!constant}

The module provides several utilities for building and combining functions:

\begin{lstlisting}[language=Lattice, caption={comp, flip, and constant}]
import "lib/fn" as F

// comp(f, g) is a module-accessible wrapper around compose()
let add_one_then_double = F.comp(|x| { x * 2 }, |x| { x + 1 })
print(add_one_then_double(4))  // 10

// flip(f) reverses the argument order of a 2-argument function
let div = |a, b| { a / b }
let inv_div = F.flip(div)
print(inv_div(2, 10))  // 5  (calls div(10, 2))

// constant(v) returns a function that always returns v
let always_zero = F.constant(0)
print(always_zero(99))     // 0
print(always_zero("hi"))   // 0
\end{lstlisting}

\subsection{apply\_n --- Repeated Application}\label{subsec:fn-apply-n}
\index{fn standard library!apply\_n}

The \lstinline|F.apply_n(f, n, value)| function applies \lstinline|f| to \lstinline|value| a total of \lstinline|n| times:

\begin{lstlisting}[language=Lattice, caption={Repeated function application}]
import "lib/fn" as F

let double = |x| { x * 2 }
print(F.apply_n(double, 0, 3))  // 3
print(F.apply_n(double, 1, 3))  // 6
print(F.apply_n(double, 4, 1))  // 16  (1 -> 2 -> 4 -> 8 -> 16)
\end{lstlisting}

This is a controlled form of recursion---useful for simulations, numerical methods, or any situation where you want to iterate a transformation a fixed number of times.

\subsection{thread() --- Module-Level Piping}\label{subsec:fn-thread}
\index{fn standard library!thread}

The \lstinline|F.thread(value, f1, f2)| function is the module's equivalent of a two-function \lstinline|pipe()|, for cases where you want to use the \lstinline|fn| module's qualified syntax:

\begin{lstlisting}[language=Lattice, caption={thread for 2-step piping}]
import "lib/fn" as F

let result = F.thread(10,
    |x| { x + 5 },
    |x| { x * 2 }
)
print(result)  // 30
\end{lstlisting}

For longer chains, the built-in \lstinline|pipe()| is more flexible since it accepts any number of functions.

\subsection{Collection Utilities}\label{subsec:fn-collections}
\index{fn standard library!group\_by}\index{fn standard library!partition}\index{fn standard library!frequencies}

Beyond sequences and composition, the \lstinline|fn| module includes utilities for working with arrays and maps.

\subsubsection{group\_by --- Categorize Elements}
\index{fn standard library!group\_by}

\begin{lstlisting}[language=Lattice, caption={Grouping elements}]
import "lib/fn" as F

let words = ["apple", "avocado", "banana", "blueberry", "cherry"]
let grouped = F.group_by(words, |w| { w[0] })

print(grouped)
// {a: [apple, avocado], b: [banana, blueberry], c: [cherry]}
\end{lstlisting}

The key function determines which group each element belongs to.
The result is a map from keys to arrays of matching elements.

\subsubsection{partition --- Split into Two Groups}
\index{fn standard library!partition}

\begin{lstlisting}[language=Lattice, caption={Partitioning an array}]
import "lib/fn" as F

let numbers = [1, 2, 3, 4, 5, 6, 7, 8, 9, 10]
let result = F.partition(numbers, |n| { n % 2 == 0 })

print(result[0])  // [2, 4, 6, 8, 10]  (matches)
print(result[1])  // [1, 3, 5, 7, 9]   (non-matches)
\end{lstlisting}

\subsubsection{frequencies --- Count Occurrences}
\index{fn standard library!frequencies}

\begin{lstlisting}[language=Lattice, caption={Counting element frequencies}]
import "lib/fn" as F

let votes = ["yes", "no", "yes", "yes", "no", "abstain"]
let tally = F.frequencies(votes)

print(tally)
// {yes: 3, no: 2, abstain: 1}
\end{lstlisting}

\subsubsection{uniq\_by --- Deduplicate by Key}
\index{fn standard library!uniq\_by}

\begin{lstlisting}[language=Lattice, caption={Deduplication by a key function}]
import "lib/fn" as F

let items = [1, 2, 3, 4, 5, 6, 7, 8, 9]
let unique_mod3 = F.uniq_by(items, |x| { x % 3 })

print(unique_mod3)  // [1, 2, 3]  (first occurrence of each mod-3 class)
\end{lstlisting}

\subsection{The Result Type}\label{subsec:fn-result}
\index{fn standard library!Result type}\index{Result type (fn module)}

The \lstinline|fn| module also provides a lightweight \lstinline|Result| type for representing success or failure without exceptions:

\begin{lstlisting}[language=Lattice, caption={The Result type from the fn module}]
import "lib/fn" as F

let success = F.ok(42)
let failure = F.err("something went wrong")

print(F.is_ok(success))   // true
print(F.is_err(failure))  // true

print(F.unwrap(success))  // 42
print(F.unwrap_or(failure, 0))  // 0

// Transform the Ok value
let doubled = F.map_result(success, |x| { x * 2 })
print(F.unwrap(doubled))  // 84

// Chain operations that might fail
let chained = F.flat_map_result(success, |x| {
    if x > 0 { F.ok(x * 10) } else { F.err("must be positive") }
})
print(F.unwrap(chained))  // 420
\end{lstlisting}

The \lstinline|F.try_fn()| wrapper catches runtime errors and returns a Result:

\begin{lstlisting}[language=Lattice, caption={Catching errors with try\_fn}]
import "lib/fn" as F

let safe = F.try_fn(|_| { 100 / 4 })
print(F.unwrap(safe))  // 25

let risky = F.try_fn(|_| { 100 / 0 })
print(F.is_err(risky))  // true
\end{lstlisting}

%% ───────────────────────────────────────────────────────────
\section{Currying and Partial Application Patterns}\label{sec:currying}
\index{currying}\index{partial application}

Currying and partial application are techniques for breaking a multi-argument function into a chain of single-argument functions.
This lets you ``pre-fill'' some arguments and pass the specialized function around.

\subsection{curry() --- One Argument at a Time}\label{subsec:curry}
\index{fn standard library!curry}

The \lstinline|F.curry()| function transforms a 2- or 3-argument function into a chain of single-argument functions:

\begin{lstlisting}[language=Lattice, caption={Currying a 2-argument function}]
import "lib/fn" as F

let add = F.curry(|a, b| { a + b })

// Now add(5) returns a new function that adds 5
let add_five = add(5)
print(add_five(3))   // 8
print(add_five(10))  // 15

// Or call it all at once via chaining
print(add(2)(3))     // 5
\end{lstlisting}

Three-argument functions work the same way:

\begin{lstlisting}[language=Lattice, caption={Currying a 3-argument function}]
import "lib/fn" as F

let volume = F.curry(|length, width, height| {
    length * width * height
})

let slab = volume(10)        // fix the length
let narrow_slab = slab(2)    // fix the width
print(narrow_slab(5))        // 100

print(volume(3)(4)(5))       // 60
\end{lstlisting}

\begin{defnbox}{Currying}
\emph{Currying} transforms a function \lstinline|f(a, b)| into \lstinline|f(a)(b)|.
Each call with one argument returns a new function that ``remembers'' the provided argument and waits for the rest.
\end{defnbox}

\subsection{partial() --- Bind Some Arguments Now}\label{subsec:partial}
\index{fn standard library!partial}

While currying requires you to provide arguments one at a time in order, \lstinline|F.partial()| lets you bind any number of leading arguments at once:

\begin{lstlisting}[language=Lattice, caption={Partial application}]
import "lib/fn" as F

let greet = |greeting, name| { greeting + ", " + name + "!" }

let hello = F.partial(greet, "Hello")
let howdy = F.partial(greet, "Howdy")

print(hello("Alice"))  // Hello, Alice!
print(howdy("Bob"))    // Howdy, Bob!
\end{lstlisting}

You can bind multiple arguments at once:

\begin{lstlisting}[language=Lattice, caption={Binding multiple arguments}]
import "lib/fn" as F

let log_message = |level, component, msg| {
    print("[" + level + "] " + component + ": " + msg)
}

let log_error = F.partial(log_message, "ERROR")
let log_db_error = F.partial(log_message, "ERROR", "database")

log_error("auth", "invalid token")
// [ERROR] auth: invalid token

log_db_error("connection timeout")
// [ERROR] database: connection timeout
\end{lstlisting}

\subsection{Currying with Iterators}\label{subsec:curry-with-iterators}

Currying and partial application become especially powerful when combined with iterators, because they let you create specialized transform functions on the fly:

\begin{lstlisting}[language=Lattice, caption={Curried functions in pipelines}]
import "lib/fn" as F

// A curried multiplier
let multiply = F.curry(|factor, x| { factor * x })

let prices = [12.50, 8.99, 24.00, 5.75]

// Create specialized multipliers
let with_tax = multiply(1.08)
let discounted = multiply(0.85)

// Use them in pipelines
let taxed_prices = iter(prices).map(with_tax).collect()
print(taxed_prices)  // [13.5, 9.7092, 25.92, 6.21]

let sale_prices = iter(prices).map(discounted).collect()
print(sale_prices)  // [10.625, 7.6415, 20.4, 4.8875]
\end{lstlisting}

Compare this with the alternative---writing a fresh closure for each operation.
Currying lets you express the intent (``multiply by 1.08'') more directly than \lstinline+|x| { x * 1.08 }+.

\begin{tipbox}{Curry or Partial?}
Use \lstinline|curry()| when you want to apply arguments one at a time and the function has 2--3 parameters.
Use \lstinline|partial()| when you want to bind one or more leading arguments in a single call, especially for functions with many parameters.
\end{tipbox}

%% ───────────────────────────────────────────────────────────
\section{When Functional Style Shines in Lattice}\label{sec:when-functional}
\index{functional programming!best practices}

Functional programming is a tool, not a religion.
Lattice gives you both imperative and functional styles; the art is knowing when each one serves you best.

\subsection{Data Transformation Pipelines}\label{subsec:when-pipelines}

Functional style is at its best when you are transforming data through a series of well-defined steps:

\begin{lstlisting}[language=Lattice, caption={A data transformation pipeline}]
import "lib/fn" as F

let raw_records = [
    "Alice,95,Math",
    "Bob,87,Science",
    "Carol,92,Math",
    "Dave,78,Science",
    "Eve,96,Math"
]

// Parse, filter, and summarize in one flow
let top_math_scores = pipe(raw_records,
    |records| { iter(records).map(|r| { r.split(",") }).collect() },
    |parsed| { iter(parsed).filter(|r| { r[2] == "Math" }).collect() },
    |math| { iter(math).map(|r| { to_int(r[1]) }).collect() },
    |scores| { iter(scores).filter(|s| { s >= 90 }).collect() }
)

print(top_math_scores)  // [95, 92, 96]
\end{lstlisting}

Each step has a clear purpose, and the whole pipeline reads like a recipe.

\subsection{Callback-Heavy APIs}\label{subsec:when-callbacks}

When working with APIs that accept callbacks---event handlers, sorting comparators, validation rules---functional tools help you build specialized callbacks without repetitive boilerplate:

\begin{lstlisting}[language=Lattice, caption={Building validators with partial application}]
import "lib/fn" as F

let is_between = |min_val, max_val, x| {
    x >= min_val && x <= max_val
}

let is_valid_age = F.partial(is_between, 0, 150)
let is_valid_score = F.partial(is_between, 0, 100)

let ages = [25, -3, 42, 200, 17]
let valid_ages = iter(ages).filter(is_valid_age).collect()
print(valid_ages)  // [25, 42, 17]
\end{lstlisting}

\subsection{Configuration and Strategy Patterns}\label{subsec:when-strategy}

Functions are values in Lattice.
You can store them in maps, pass them as arguments, and return them from other functions.
This makes the strategy pattern natural:

\begin{lstlisting}[language=Lattice, caption={Strategy pattern with functions}]
let strategies = Map::new()
strategies.set("uppercase", |s| { s.upper() })
strategies.set("lowercase", |s| { s.lower() })
strategies.set("reverse", |s| { s.reverse() })

fn apply_strategy(strategy_name: String, text: String) -> String {
    let transform = strategies.get(strategy_name)
    return transform(text)
}

print(apply_strategy("uppercase", "hello"))   // HELLO
print(apply_strategy("reverse", "Lattice"))   // ecittaL
\end{lstlisting}

\subsection{When Imperative Style Is Better}\label{subsec:when-imperative}

Functional style is not always the clearest choice.
Prefer imperative code when:

\begin{itemize}[nosep]
  \item The logic involves complex branching or early returns that don't fit neatly into a pipeline.
  \item You need to accumulate state across multiple collections simultaneously.
  \item The code would require deeply nested function calls that obscure the intent.
  \item Performance matters and you need fine-grained control over allocation.
\end{itemize}

\begin{lstlisting}[language=Lattice, caption={Sometimes a loop is clearer}]
// Functional version -- hard to follow
let result = pipe(items,
    |xs| { iter(xs).filter(|x| { x > threshold }).collect() },
    |xs| {
        iter(xs).reduce(|acc, x| {
            if x > acc { x } else { acc }
        }, 0)
    }
)

// Imperative version -- immediately clear
flux max_val = 0
for item in items {
    if item > threshold && item > max_val {
        max_val = item
    }
}
\end{lstlisting}

The best Lattice code uses both styles where each is strongest.
Functional pipelines for data transformations; imperative loops for stateful logic.

\begin{tipbox}{The Litmus Test}
If you have to read a functional pipeline more than twice to understand what it does, it is probably trying to be too clever.
Break it into named intermediate variables or switch to a loop.
Clarity always beats cleverness.
\end{tipbox}

%% ───────────────────────────────────────────────────────────
\section{Putting It All Together}\label{sec:fp-putting-together}

Let's close with a larger example that combines iterators, \lstinline|pipe()|, currying, and the \lstinline|fn| module into a cohesive program.

\begin{lstlisting}[language=Lattice, caption={A complete functional example}]
import "lib/fn" as F

// --- Data ---
let inventory = [
    ["Widget A", 29.99, 150],
    ["Widget B", 9.99, 500],
    ["Gadget C", 49.95, 30],
    ["Gadget D", 14.50, 200],
    ["Doohickey E", 4.99, 1000]
]

// --- Helpers built with currying ---
let min_price = F.curry(|threshold, item| { item[1] >= threshold })
let min_stock = F.curry(|threshold, item| { item[2] >= threshold })

// --- Pipeline: find valuable items with good stock ---
let report = pipe(inventory,
    |items| { iter(items).filter(min_price(10.0)).collect() },
    |items| { iter(items).filter(min_stock(100)).collect() },
    |items| {
        iter(items).map(|item| {
            let name = item[0]
            let total_value = item[1] * item[2]
            name + ": $" + to_string(total_value)
        }).collect()
    }
)

for line in report {
    print(line)
}
// Widget A: $4498.5
// Widget B: $4995.0
// Gadget D: $2900.0
\end{lstlisting}

Each piece is small and testable on its own.
The curried predicates can be reused across different pipelines.
The \lstinline|pipe()| call makes the overall flow readable.
And because the iterator operations are lazy where possible, no unnecessary work is done.

%% ───────────────────────────────────────────────────────────
\section{Exercises}\label{sec:fp-exercises}

\begin{enumerate}
  \item \textbf{Compose Three.}
    Lattice's built-in \lstinline|compose()| takes exactly two functions.
    Write a function \lstinline|compose3(f, g, h)| that composes three: \lstinline|compose3(f, g, h)(x) = f(g(h(x)))|.
    Test it by composing three arithmetic operations.

  \item \textbf{Curried Filter.}
    Using \lstinline|F.curry()|, create a reusable \lstinline|longer_than| predicate that filters strings by minimum length.
    Use it to extract all words longer than 5 characters from the sentence \lstinline|"The quick brown fox jumped over the lazy dog"|.

  \item \textbf{Word Frequency Pipeline.}
    Given a block of text (a multi-line string), use \lstinline|pipe()| and the \lstinline|fn| module to:
    \begin{enumerate}[nosep]
      \item Split the text into words.
      \item Convert all words to lowercase.
      \item Count the frequency of each word using \lstinline|F.frequencies()|.
    \end{enumerate}
    Print the resulting frequency map.

  \item \textbf{Result Pipeline.}
    Write a sequence of operations using \lstinline|F.ok()|, \lstinline|F.flat_map_result()|, and \lstinline|F.map_result()| to:
    \begin{enumerate}[nosep]
      \item Start with a string that might represent a number.
      \item Parse it to an integer (returning \lstinline|F.err()| if it fails).
      \item Double the number.
      \item Check that it is positive (returning \lstinline|F.err()| if not).
    \end{enumerate}
    Test with both valid and invalid inputs.
\end{enumerate}

%% ───────────────────────────────────────────────────────────
\bigskip
\begin{center}\rule{0.5\linewidth}{0.4pt}\end{center}
\medskip

\noindent\textbf{What's Next}\quad
With iterators and functional tools in your kit, you are well-equipped to write expressive, composable Lattice code.
In \cref{ch:modules-imports}, we turn to a different kind of composition: organizing your code into \emph{modules} and \emph{packages}, so that the functions and types you build can be shared, reused, and managed as your projects grow.
